\scanpage{73}
\mspage{- 18 -}

\[
f_i(x_1,..,x_p\,;\,v_1(x_1,..,x_p),..,v_q(x_1,..,x_p))=0
\quad(1\leq i\leq q)\tag*{(2)}
\]

\uline{dans un voisinage de} $(a_1,..,a_p)$, \uline{on a} $u_i=v_i$ $(1\leq i\leq q)$ \uline{dans un voisinage de} $(a_1,..,a_p)$.

Par tra\handwritten{\unsure{-|{}}}nslation, on peut se ramener au cas où $a_j=0$ pour $1\leq j\leq p$ et $b_i=0$ pour $1\leq i\leq q$, . Posons pour $1\leq i\leq q$

\[
f_i(x_1,..,x_p;y_1,..,y_q)=f_{i0}+
\handwritten{\sum_{j=1}^{q}}f_{ij}y_j+
\handwritten{\sum_{(n_j)}}f_{\unsure{1|i};n_1..n_q}y_1^{n_1}\ldots y_q^{n_q}\tag*{(3)}
\]

les $f_{i0}$, $f_{ij}$ et $f_{i;n_1..n_q}$ étant des séries entières par rapport aux $x_k$ $(1\leq k\leq p)$ sommables dans un polycylindre ouvert $P$ de centre l'origine dans $K^p$, $f_{i0}$ n'ayant pas de terme constant ; en outre, $f_i$ est sommable dans $P\times Q$, où $Q$ est un polycylindre de centre l'origine dans $K^q$, et dans la seconde sommation des seconds membres de (3), on a toujours $\handwritten{\sum_{j=1}^{q}}n_j\geq 2$. Par hypothèse, la matrice $\underline F=(f_{ij})$ est holomorphe dans $P$ et inversible à l'origine donc elle est inversible dans un polycylindre ouvert contenu dans $P$, et son inverse $\underline G=(g_{ij})$ est holomorphe dans ce polycylindre ; nous supposerons donc chacune des $g_{ij}$ développées en série entière sommable dans un polycylindre ouvert $P_1\handwritten{\subset} P$ de centre l'origine ; si on pose alors $g_i=\handwritten{\sum_{j=1}^{q}}g_{ij}f_j$, on aura

\[
g_i(x_1,..,x_p;y_1,..,y_q)=-h_{i0}+y_i-
\handwritten{\sum_{(n_j)}}h_{i;n_1..n_q}y_1^{n_1}\ldots y_q^{n_q}\tag*{(4)}
\]

où $h_{i0}$ et les $h_{i;n_1..n_q}$ sont des séries entières par rapport aux $x_k$, sommables dans $P_1$, $h_{i0}$ n'ayant pas de terme constant, et la somme étant étendue aux $(n_j)$ tels que $\handwritten{\sum_{j=1}^{q}} n_j\geq 2$. La première partie du théorème revient à prouver qu'il existe $q$ séries entières sans terme constant $u_i$ $(1\leq i\leq q)$ par rapport aux $x_k$, sommables dans un voisinage de l'origine, et telles que l'on ait

\[
u_i=h_{i0}+\handwritten{\sum_{(n_j)}}h_{i;n_1..n_q}u_1^{n_1}\ldots u_q^{n_q}
\qquad\text{pour }\ 1\leq i\leq q\ .\quad .\tag*{(5)}
\]
% archive sheet number: 72

